\documentclass[10pt,a4paper]{amsart}
\usepackage[margin=19mm]{geometry}
\usepackage{amsmath,amssymb,mathtools}
\usepackage[scr=rsfs,cal=euler]{mathalfa}
\usepackage{xfrac}
\usepackage[unicode,hidelinks,bookmarks=false]{hyperref}
\newcommand{\ie}{i.e.\ }
\newcommand{\cf}{cf.\ }
\newcommand{\resp}{respectively\ }
\newcommand{\Subsection}{\subsection}
\providecommand{\nobreakdash}{\nobreak\hbox{-}}
\setlength{\emergencystretch}{2em}
\multlinegap0pt
\usepackage{amssymb}
\usepackage{mathtools}
\usepackage[shortlabels]{enumitem}
\newtheorem{theorem}{Theorem}
\newtheorem{lemma}{Lemma}
\usepackage{cleveref}
\crefname{theorem}{Theorem}{Theorems}
\crefname{lemma}{Lemma}{Lemmas}
\title{Conditional upper bounds on the least character non-residue}
\author{Aritro Pathak}
\date{Source: arXiv:2405.13316v2 (CC BY 4.0)}
\begin{document}
\maketitle

\noindent\emph{Source and licence.} Conditional upper bounds on the least character non-residue. Version arXiv:2405.13316v2 (CC BY 4.0), distributed by arXiv under the Creative Commons Attribution 4.0 International licence, http://creativecommons.org/licenses/by/4.0/, as recorded in the arXiv licence metadata for this exact version and shown on the abstract page https://arxiv.org/abs/2405.13316v2. Full licence notice: \texttt{licenses/arxiv-2405.13316-CC-BY-4.0.txt}.\par\smallskip

\noindent\textit{Editorial scope.} Counted unit: the article's theorem mainthm (source0.tex lines 163--180). The article's complete original proof of it is delivered with it (source0.tex lines 269--411, one \texttt{proof} environment of 143 lines, which the article places away from the statement). No mathematics is written, repaired or re-derived: the statement and the proof are the article's own lines, in the article's own notation, and no retained line is substituted or re-flowed.  Retained uncounted as the source-local context the proof needs: lines 199--209.  Nothing was omitted from any retained span, so the package carries no omission record.  The retained lines cite 1 works, whose entries are transcribed verbatim from the article's own bibliography source in the e-print and delivered with short editorial labels recorded in the item's reference mapping.  No retained line contains a figure environment, an image-inclusion command or drawing code, so no image is delivered and the package declares no asset dependency.  Editorial formatting only, in addition: an AMS \texttt{amsart} preamble that restates the article's own declarations and macros used by the retained lines, namely the environments \texttt{theorem}, \texttt{lemma} on separate counters, together with the standard packages those lines need.  The retained environments print in this document as Theorem 1, Lemma 1 in order of appearance, whereas the article prints the same items with the numbering of its own full document.  The complete native source of the article as distributed in the arXiv e-print for this exact version is bundled unchanged as \texttt{sources/159519/source0.tex}.
\begin{lemma}\label{lemma}
    Let $n(r;t,\chi)$ denote the number of zeros $\rho$ of $L(s,\chi)$ in the disc 
$|\rho-(1+it)|\le r$. Then
\[
n(r;t,\chi)\ll r\log(q\tau)
\]
uniformly for
\[
\frac{1}{\log(q\tau)} \le r \le \frac{3}{4}.
\]
\end{lemma}

\begin{theorem}\label{mainthm}

There exist absolute constants $C>0$ and $K_1>0$ such that the following holds.
Let $|t_0|>1$, let $\chi$ be a non-principal character $(\bmod\, q)$ with
$q \ge e^{|t_0|}+4$, and let
\[
\frac{1}{\log(|t_0|+4)} \le \delta \le \frac{1}{2}.
\]
If
\[
L(s,\chi)\neq 0 \quad \text{for } 1-\delta<\sigma<1, \qquad
|t| \le K_1 |t_0|^2 \log q,
\]
then
\[
n(\chi) < C\left( K_1 |t_0|^2 \frac{\log q}{\delta} \right)^{1/\delta}.
\]
\end{theorem}

\begin{proof}[Proof of \cref{mainthm}]
 Let $x(q,t_0)=\Big(\frac{K_2 t_{0}^{2}\log q}{\delta}\Big)^{\frac{1}{\delta}},y(q,t_0)=\sqrt{C}$  for constants $K_2,C$ , such that $2\leq y(q,t_0)\leq x(q,t_0)^{1/3}$. Henceforth, the dependence on $q, t_0$ of these $x,y$ parameters should be understood, and we don't write them as functions of $q,t_0$ explicitly.

 \bigskip
 
As will be seen in the course of the proof, the kernel function that works for our purpose, is
\begin{equation}
    K(s,t_0)=\int_{\frac{x}{y^{2}}}^{xy^{2}} \Big(2 \log y -\Big|\log \frac{x}{u}\Big|\Big)\Big(\frac{x}{u}\Big)^{1+it_0}u^{s-1} du,
\end{equation}
where the weight function is $w(u)=\big(2 \log y -|\log \frac{x}{u}|\big)\big(\frac{x}{u}\big)^{1+it_0}$

%\footnote{This reduces to the weight function used in the proof of Theorem 13.12 in \cite{Mont} when $t_0 =0$.}.

This weight function is  non-zero within the interval $(x/y^{2}, xy^{2})$ and take $w(u)=0$ elsewhere.

With a change of variable: $\frac{x}{u}=m$, this integral becomes 
\begin{multline}
    K(s,t_0)=x^{s}\int_{\frac{1}{y^{2}}}^{y^{2}} \Big( 2\log y -|\log m|  \Big) m^{it_0 -s} dm=2\log y\cdot x^{s} \Big( \frac{m^{it_0 -s+1}}{it_0 -s+1} \Big)|_{\frac{1}{y^{2}}}^{y^{2}}\\
    +x^{s}\int_{1/y^{2}}^{1} (\log m) m^{it_0 -s}dm -x^{s}\int_{1}^{y^{2}}(\log m) m^{it_0 -s} dm.
\end{multline}

\noindent This gives us:

\begin{multline}
    K(s,t_0)=\frac{2(\log y) x^{s} y^{2(it_0 -s +1)} }{it_0 -s+1}-\frac{2(\log y) x^{s} y^{-2(it_0 -s +1)} }{it_0 -s+1} -\frac{2x^{s}}{(it_0 -s+1)^{2}}\\ +\frac{2(\log y) x^{s} y^{-2(it_0 -s +1)} }{it_0 -s+1}-\frac{2(\log y) x^{s} y^{2(it_0 -s +1)} }{it_0 -s+1}+\frac{x^{s}y^{-2(it_0 -s+1)}}{(it_0-s+1)^{2}}+\frac{x^{s}y^{2(it_0 -s+1)}}{(it_0-s+1)^{2}}.
\end{multline}

Upon cancellations among four terms above, we get, 
\begin{align}
    K(s,t_0)=\frac{-2x^s +x^s y^{-2(it_0 -s+1)} +x^s y^{2(it_0 -s +1)} }{(it_0 -s+1)^2}.
\end{align}

\noindent Thus, we have for the kernel $K(s,t_0)$,

\begin{equation}
K(s,t_0)=x^{s}\Big(\frac{y^{s-it_0 -1}- y^{1+it_0 -s}}{s-it_0 -1}\Big)^{2}.
\end{equation}

%After this, it can be shown that the estimates from Equation 11.5 and 10.17 work for any general $t$, as it does for the the case of $t=0$, and the later estimates of the corresponding integrals will also work similar to the case of $t=0$.

\bigskip 

 Let $n(r;t,\chi)$ denote the number of zeros of $L(s,\chi)$ in the disk $D= \{x\in \mathbb{C}:|x-(1+it)|<r\}$. For any $t_0 \in \mathbb{R}$ with $\frac{1}{\log q(|t_0|+4)}\leq R \leq 1$, we have the estimate,
\begin{equation}\label{impeq}
    \sum\limits_{|\rho -1-it_0|>R} \frac{1}{|\rho- 1-it_0|^{2}} \ll \frac{\log q(|t_0|+4)}{R}.
\end{equation}

\noindent This follows from the fact that 

\begin{equation}
    \sum\limits_{R<|\rho -1-it_0|<2R} \frac{1}{|\rho -1-it_0|^{2}} \ll \frac{1}{R^{2}} n(2R;t_0,\chi) \ll \frac{\log q(|t_0|+4)}{R},
\end{equation}
where the first inequality follows because each individual summand on the left is bounded by $1/R^{2}$ and the number of terms is bounded by $n(2R;t_0,\chi)$. From the estimates on $n(R;t_0,\chi)$ from \cref{lemma}, %\footnote{See for example Theorem 11.5 in \cite{Mont}. While there is an upper bound of $3/4$ on $r$ in the estimate on $n(r,t,\chi)$, for the case where $3/4 \leq r\leq 1$ one can further use an additional factor of $\log(q(|t_0|+4))$ for the number of zeros in the region  $0\leq \beta \leq 1$ and $t_0 \leq \gamma \leq t_0 +1$, akin to the estimate used in the ensuing paragraph.} 
we have the second inequality. Now replacing $R$ by $2^{k}R$, and summing over sufficiently many terms, we have a count of all the zeros within the annulus $R<|\rho -1-it_0|<1$, and we incorporate a constant pre-factor arising from the finite geometric progression, to get,
\begin{equation}\label{eq16}
    \sum\limits_{R<|\rho -1-it_0|<1} \frac{1}{
    |\rho -1-it_0|^{2}}\ll \frac{\log q(|t_0|+4)}{R}.
\end{equation}

For zeros further away from the point $(1+it_0)$, we use the estimate that the number of zeros of the function $L(s,\chi)$ within the range $0\leq \beta \leq 1$ and $T\leq \gamma \leq T+1$ is $\ll \log (q(|T|+4))$. Therefore putting $K=1$ in \cref{otherimp} below, it can be seen using \cref{eq16} along with an easy comparison that the bound in \cref{impeq} is attained. 

The zero-free region in the hypothesis of the theorem is symmetric with respect to the line $y=t_0$, and it is clear that for any positive integer $K$, we attain an upper bound of the form,

\begin{multline}\label{otherimp}
     \sum\limits_{|\rho-1-it_0|>K} \frac{1}{|\rho-1-it_0|^{2}} \ll \sum\limits_{n=K}^{\infty} \frac{\log q(n+|t_0|+4)}{n^{2}} \\ \ll \frac{\log q}{K}+\frac{\log(K+|t_0|+4)}{K}+ \Big(\frac{1}{|t_0|+4}\Big)\log\Big(1+ \frac{|t_0|+4}{K}\Big).
\end{multline}

%Further, following Theorem 11.5 in \cite{Mont}, we also have 



The weight function being $w(u)= (2\log y -|\log(x/u)|)(x/u)^{1+it_0}$ in the interval $2\leq y \leq x^{1/3}$, we get for the weighted sum\footnote{See for example, the corresponding equation 13.31 in \cite{Mont}.},


\begin{equation}\label{eq:mainone}
\sum\limits_{n} w(n)\chi(n)\Lambda(n)= \frac{-1}{2\pi i} \int_{\sigma_0 -i\infty}^{\sigma_0 +i\infty} \frac{L'}{L}(s,\chi) \Bigg(\frac{y^{s-1-it_0}-y^{1+it_0-s}}{s-it_0-1} \Bigg)^{2} x^{s}ds.
\end{equation}


\noindent Here we have taken $\sigma_0 >1$ since $L(s,\chi)$ may have zeros in the critical strip and thus $\frac{L'(s,\chi)}{L(s,\chi)}$ would have simple poles in the critical strip. We move the contour to the left of the imaginary axis, to the line $\sigma_0 = -1+\theta$ for some small $\theta$ and get from \cref{eq:mainone},

\begin{multline}\label{eq:imp2}
    -\sum\limits_{\rho}\Big( \frac{y^{\rho -1-it_0}-y^{1+it_0-\rho}}{\rho -1-it_0}  \Big)^{2}x^{\rho} -\frac{(1-\kappa)}{(1+it_0)^{2}}(y^{1+it_0}- 1/y^{1+it_0})^{2} \\ -\frac{1}{2\pi i} \int_{-1+\theta -i\infty}^{-1+\theta +i\infty} \frac{L'}{L}(s,\chi) \Big( \frac{y^{s-1-it_0}-y^{1+it_0-s}}{s-it_0-1} \Big)^{2}x^{s}ds.
\end{multline}

\noindent The second term arises because $L(s,\chi)$ has a zero at $s=0$ when $\chi(-1)=1$ in which case $\kappa=0$ and otherwise $\kappa=1$.

If $\chi$ is induced by a primitive character $\chi^{*}$, then from the standard expression for $L(s,\chi)$ in terms of $L(s,\chi^{*})$, we take the logarithmic derivative and get that,

\begin{equation*}
    \frac{L'}{L}(s,\chi)=\frac{L'}{L}(s,\chi^{*})+\sum\limits_{p|q} \frac{\chi^{*}(p)\log p}{p^{s}-\chi^{*}(p)}.
\end{equation*}

In the sum above, we have $\sigma=-1+\theta$, and so it is easily seen that the summand is $\ll \log p$ and thus the sum of $\log p$ over all the primes $p|q$ is trivially bounded by $\log q$ itself. Further, since we have chosen the contour line to be where $\sigma_0=-1+\theta$, by another standard result (See for example Theorem 12.9 of \cite{Mont}), we can bound the above logarithmic derivative term for the primitive character by $\ll \log 2q(|t|+4)$, and then easily get a finite contribution from the contour integral and conclude that the total contribution to \cref{eq:imp2} from these terms is bounded by
\begin{equation*}
\ll x^{-1+\theta}y^{4-2\theta} \log q(|t_0|+4)\leq x^{-1+\theta}y^{4} \log q(|t_0|+4).
\end{equation*}

\noindent This will turn out to be smaller than the contribution from the first sum in \cref{eq:imp2}: recall that we started out with the bound $2\leq y \leq x^{1/3}$, and we will consider $y=O(1)$.  

Now first consider the contribution to the sum in \cref{eq:imp2} from the zeros $\rho=\beta +i\gamma$, with $\beta \leq 1-\delta$. In this case, considering $R=\delta$ in \cref{impeq}, we get the bound,

$$\ll \sum\limits_{|\rho-1-it_0|>\delta} \frac{x^{1-\delta}(3+y^{2})}{|\rho -1-it_0|^{2}} \ll x^{1-\delta}(3+y^{2}) \frac{(\log q(|t_0|+4))}{\delta}.$$

\noindent Further if $\rho$ is a zero in the critical strip with $\beta >1-\delta$, then by hypothesis, we have $|\gamma-t_0|\geq K(q,t_0)=K_{1}|t_0|^{2}\log q$. In order to use \cref{otherimp}, one can choose $K_1$ large enough so that $ K=\lfloor K(q,t_0) \rfloor$ is a positive integer. Each of these summands in \cref{eq:imp2} is bounded from above by $x(3+y^{2\delta})/(|\rho-1-it_0|^{2})$, when $1-\delta<\beta<1$.

%Now if $\delta^{2}\log q(|t_0|+4)<1$ then we use the bound from \cref{impeq} with $R=\delta^{2}\log q(|t_0|+4)$ to conclude that the total contribution of these zeros is bounded by $ \ll x(3+y^{2\delta})/\delta^{2}$, since



We see that when $x$ is large enough, both of these bounds are greater than the bound of $x^{-1+\theta}y^{4}\log q(|t_0|+4)$ obtained from the contour integral in \cref{eq:imp2}. Combining with the estimate from \cref{otherimp}, we have that

%\footnote{Even in a thin rectangle where $1-\delta<\sigma<1$ and $K \leq t\leq K+1$ with $\delta$ arbitrarily small, one only gets an upper bound of the form $\ll\log q(K+4)$ on the number of zeros in the rectangle, and the implied constant is independent of $\delta$. This follows from Theorem 11.5 of \cite{Mont} by writing this rectangle as a union of $1/\delta$ rectangles where in each such rectangle we have a bound of the form $\delta\log q(K+4)$.}

\begin{multline}
     |\sum_{n}w(n)\chi(n)\Lambda(n)|\leq c\Bigg(\Big(x^{1-\delta}(3+y^{2}) \frac{(\log q(|t_0|+4))}{\delta}\Big)\\ + x(3+ y^{2\delta})\Bigg(\frac{\log q}{K}+ \Big(\frac{1}{|t_0|+4}\Big)\log\Big(1+ \frac{|t_0|+4}{K}\Big)+\frac{\log(K+|t_0|)}{K} \Bigg)\Bigg).
\end{multline}

\noindent If we now choose $x$ and  $K_1$ above so that $K=\delta x^{\delta}$ , then we will have $x(q,t_0)=\Big(\frac{K_2 t_{0}^{2}\log q}{\delta}\Big)^{\frac{1}{\delta}}$ for some constant $K_2$ as written in the beginning of the proof. We obtain,
\begin{multline}\label{eq:oneimp}
     |\sum_{n}w(n)\chi(n)\Lambda(n)|\leq c'\Big(x^{1-\delta}(3+y^{2}) \frac{(\log q(|t_0|+4))}{\delta}+\frac{1}{\delta}x^{1-\delta}(3+y^{2\delta})\big(\log (\delta x^{\delta}+|t_0|)+\log q \big)\\ +\Big(\frac{1}{|t_0|+4}\Big)\log\Big(1+ \frac{|t_0|+4}{\delta x^{\delta}}\Big) x(3+y^{2\delta})\Big).
\end{multline}

Next we estimate the sum  $\sum\limits_{n}w(n)\chi_0(n)\Lambda(n)$ where $\chi_0$ is the principal character modulo $q$. In that case, from \cref{eq:mainone}, putting in the principal character, we would then proceed as in the proof of the Prime Number Theorem, and then the leading order term would come from the residue of the function $-\frac{L'}{L}(s,\chi_0) \Big(\frac{y^{s-1-it_0}-y^{1+it_0-s}}{s-it_0-1} \Big)^{2}x^{s}$. We have a simple pole at $s=1$ and none at $s=1+it_0$. We have the residue at $s=1$ 
%\footnote{We note from the behavior of the $sinc$ function that the expression has the maximum value of $4x (\log y)^{2}$ when $t_0=0$ and for larger absolute values of $t_0$, we look at the `peaks' where the sine values are $\pm 1$. },

\begin{equation}\label{eq:last}
    -\frac{x}{t_{0}^{2}}\big(y^{-it_0} -y^{it_0}\big)^{2}=\frac{4x}{t_{0}^{2}}\sin^{2}(t_0 \log y).
\end{equation}
For the $y$ values we restrict to the sequence of values of $y_k$ so that $t_0 \log y_k =\frac{\pi}{2}(2k+1)$ (as well as $y_k\geq 2$) , i.e., we choose the values $y_k=e^{\frac{\pi}{2t_0}(2k+1)}$. Thus, $\sin^{2}(t_0 \log y_k)=1$ and then the leading order term in the sum $\sum\limits_{n}w(n)\chi_0(n)\Lambda(n)$ is $\frac{4x}{t_{0}^{2}}$.  Note that given any value of $t_0$ with $|t_0|>1$, we can choose $k$ appropriately so that $y(q,t_0)=O(1)$ and $y(q,t_0)\geq 2$.

We take $x$ large enough so that
\begin{equation}\label{eq24}
\delta x^{\delta}\gg (|t_0|+4),
\end{equation}
in which case the third term on the right of \cref{eq:oneimp} also has a factor of $x^{1-\delta}$, upon using the approximation of $\log\big(1+ \frac{|t_0|+4}{\delta x^{\delta}}\big)$by $ \frac{|t_0|+4}{\delta x^{\delta}} $. Given $|t_0|>1$, we consider $q\geq e^{|t_0|+4}$, in which case $\log q(|t_0|+4)$ is approximated by $ \log q$. In this case, noting that $y(q,t_0)=O(1)$, the dominant term on the right of \cref{eq:oneimp} is $c x^{1-\delta}\frac{\log q}{\delta}$.

If we require for $q\geq e^{|t_0|+4}$ , and with some large enough absolute constant $K_1$ that
\begin{align}
    \frac{4x}{t_{0}^{2}}\geq K_1 x^{1-\delta}\frac{\log q}{\delta}\Leftrightarrow 4x^{\delta}\delta \geq K_1 t_{0}^{2}\log q,
\end{align}
in which case \cref{eq24} is also satisfied, then the character $\chi(\text{mod} \ q)$ can't equal the principal character $\chi_0(\text{mod} \ q)$ till the value of $xy^{2}=Cx=C(\frac{K_2 t_{0}^{2}\log q}{\delta})^{\frac{1}{\delta}}$, due to the equality in Eq. 24 and the inequality in Eq. 23, and in which case we have $n(\chi)<C(\frac{K_2 t_{0}^{2}\log q}{\delta})^{\frac{1}{\delta}}$.
\end{proof}

\begin{thebibliography}{99}
\bibitem[Mont]{Mont} Montgomery, Hugh L., and Robert C. Vaughan. Multiplicative number theory I: Classical theory. No. 97. Cambridge Studies in Advanced Mathematics, Cambridge university press, 2007.
\end{thebibliography}
\end{document}
